\section{Introduction}
\label{sec:intro}

A store that records what an organization believes about its entities is asked
two kinds of temporal question. One is about the world: who held a position in
1998, which services a system depended on last March. The other is about the
record: what the store would have answered on a given day, before a correction
arrived. Temporal databases name the two axes valid time and transaction
time~\cite{snodgrass1985taxonomy,jensen1998glossary}, and a store that keeps
both is called bitemporal. A graph written by upsert keeps neither. It answers
``this entity has no owner'' in the same way whether nobody said so, somebody
said so and withdrew it, or two sources disagreed and the later write won, and
it cannot be asked which afterwards, because the rows that would distinguish
the cases were never written.

Hanzo cloud's \texttt{/v1/graph} keeps every assertion and derives answers from
them. Its store is one append-only table. A row states that an entity stands in
a relation to a value from an instant, optionally until a second instant, with
the source, the evidence and the filing identity beside it. There is no node
table, no update statement and no delete outside erasure; a retraction is a
row. Reads resolve the rows about one (entity, relation) pair at a point with
two coordinates: the valid instant asked about, and the transaction instant up
to which the store's knowledge is admitted. The design is specified in
HIP-1198~\cite{hip1198} and implemented in \texttt{apps/graph} of Hanzo
cloud~\cite{hanzocloud}, over an algebra in the package \texttt{claim} that the
ground-truth capability~\cite{hip1261} shares.

This paper describes that model as it stands at commit \texttt{1fffdf408} on
the main branch of \texttt{github.com/hanzo-inc/cloud}, states the properties
that follow from its definitions, and reports two measurements. The first is of
the store itself on CronQuestions~\cite{saxena2021cronquestions}, a temporal
question-answering benchmark over Wikidata facts with start and end years. The
second, on LongMemEval-S~\cite{wu2025longmemeval}, is of a retrieval engine in
the same memory stack and is reported for what it measures, which is not this
store.

\paragraph{What the measurements found.}
A lane in \texttt{hanzoai/benchmarks}~\cite{hanzobenchmarks} turns each
CronQuestions test question's annotation into store calls and answers from the
returned intervals with a rule shared by every store. On the
\res{cron/n} covered questions, a control that answers the same calls from
dictionaries scores \res{cron/kg/all}\% Hits@1, which is the ceiling of the
plan. Semantica~0.7.0~\cite{semantica} and the store both reach it and give the
control's answer to every question \meas{}. The store answers the median call in
\res{cron/wire/p50}\,ms and semantica in \res{cron/sem/p50}\,ms.

The first run of the same lane, on the version of the store that preceded
\texttt{1fffdf408}, scored \res{cron/first/wire/all}\% and answered none of the
\res{cron/n/se} point-in-time questions \meas{}. That version had one read
instant, \texttt{as\_of}, and it bounded the knowable instant, which is
transaction time: a fact about 1998 filed today was invisible to a question
about 1998. With the write clock replayed to each fact's own instant, the same
version scored \res{cron/first/replay/all}\%; the remaining misses came from a
second rule, under which a retraction ended every value of a pair, so the end of
one holder's term ended every other holder's. Section~\ref{sec:time} describes
the two-instant read that replaced both rules, and \S\ref{sec:eval} reports
the two defects the harness found in the replacement before it was merged.

\paragraph{Status of every claim.}
A sentence marked \meas{} states a figure from a run named in
\texttt{results-data.tex}, which is read from the artifacts the lane committed.
A sentence marked \shipped{} describes the code at \texttt{1fffdf408}. A sentence
marked \proposed{} describes a design that has not been built. Propositions are
proved from the definitions in \S\ref{sec:model}; they describe the rule, and
the code is claimed to implement the rule only as far as its tests check it.

\paragraph{Contributions.}
\begin{itemize}[leftmargin=*]
\item A description of an append-only assertion store in which each statement
  carries its own valid interval, the transaction instant is derived by the
  server rather than supplied by the filer, versions of a statement supersede
  each other only by being learned later, and a relation's cardinality is
  declared (\S\ref{sec:model}).
\item The resolution rule for a point with a valid and a transaction
  coordinate, with the properties that follow from it: stability of an answer
  as known at an instant, refusal of backdating, independence from storage
  order, and the condition under which a valid-time answer does not depend on
  when facts were filed (\S\ref{sec:model}, \S\ref{sec:time}).
\item Schema declared as assertions and checked as known when the checked fact
  became knowable (\S\ref{sec:schema}); diff, shortest path, community detection
  and a citation-checked answer built on the same rule (\S\ref{sec:queries}); and
  erasure with a receipt, with what the receipt does and does not establish
  (\S\ref{sec:erasure}).
\item A measurement on CronQuestions beside a dictionary control and an
  interval-based temporal store, including the failed first run and the defects
  found during the correction (\S\ref{sec:eval}).
\end{itemize}
